\documentclass[aspectratio=169,10pt,english]{beamer}

\input{../../../_preambulo_beamer.tex}
\input{../../../_comandos_beamer.tex}

\selectlanguage{english}

\title{MA1001B - Statistical Modeling for Decision Making}
\subtitle{Lesson 27: Z Confidence Interval for a Mean (sigma known)}
\author{}
\institute{Tecnol\'ogico de Monterrey}
\date{}

\begin{document}

\begin{frame}[plain]
  \vspace{1.2cm}
  {\Large\bfseries MA1001B - Statistical Modeling for Decision Making\par}
  \vspace{0.7cm}
  {\large Lesson 27: Z Confidence Interval for a Mean ($\sigma$ known)\par}
  \vspace{0.7cm}
  {\color{TecRojo}\rule{\textwidth}{1.2pt}}
  \vspace{0.8cm}

  MA1001B Faculty\\[0.3cm]
  Term\\[0.3cm]
  Tecnol\'ogico de Monterrey
\end{frame}

\begin{frame}{The Known-Sigma Interval Question}
  \begin{alertblock}{Guiding question}
    How is a 95\% interval for an unknown mean built when $\sigma$ is treated
    as known?
  \end{alertblock}

  \vspace{0.6em}
  By the end of this lesson, you can:
  \begin{itemize}
    \item compute $\mathrm{SE}=\sigma/\sqrt{n}$ from a known $\sigma$;
    \item form $\bar{x}\pm 1.96\cdot\sigma/\sqrt{n}$;
    \item interpret the interval as a range for $\mu$, not for one bottle;
    \item explain why unknown $\sigma$ makes this z procedure the wrong tool.
  \end{itemize}

  \vfill
  \scriptsize All fill volumes used today are synthetic. No real company data are used.
\end{frame}

\begin{frame}{A Synthetic Filling Line with Known $\sigma$}
  \small
  \begin{center}
    \begin{tabular}{lr}
      \toprule
      \textbf{Quantity} & \textbf{Value} \\
      \midrule
      Sample size $n$ & 36 bottles \\
      Known $\sigma$ & 10 ml \\
      Hidden $\mu$ (not used by the analyst) & 502 ml \\
      Point estimate $\bar{x}$ (seed 42) & 502.738498 ml \\
      \bottomrule
    \end{tabular}
  \end{center}

  \vspace{0.5em}
  \begin{exampleblock}{Standard error}
    \[
      \mathrm{SE}=\frac{\sigma}{\sqrt{n}}=\frac{10}{\sqrt{36}}=1.666667.
    \]
    The analyst knows $\sigma$ from long-run process capability and does not
    know $\mu$.
  \end{exampleblock}
\end{frame}

\begin{frame}[fragile]{Step 1: Sample Mean and Standard Error}
  \begin{columns}[T]
    \begin{column}{0.42\textwidth}
      \small
      \begin{lstlisting}
se = sigma / np.sqrt(n)
xbar = np.mean(sample)
      \end{lstlisting}

      \begin{block}{Verified output}
        $n=36$, $\sigma=10$\\
        $\bar{x}=502.738498$ ml\\
        $\mathrm{SE}=1.666667$ ml
      \end{block}
    \end{column}
    \begin{column}{0.58\textwidth}
      \centering
      \includegraphics[width=\textwidth]{../figures/z_ci_mean_01_sample_standard_error.png}
      \vspace{0.2em}
      \scriptsize Histogram and SE band from Step 1
    \end{column}
  \end{columns}
\end{frame}

\begin{frame}{Multiply $z^{*}$ by the Standard Error}
  \small
  Textbook 95\% critical value:
  \[
    z^{*}=1.96.
  \]
  Margin of error and interval:
  \[
    E=1.96\times 1.666667=3.266667,
  \]
  \[
    \bar{x}\pm E=[499.471831, 506.005165].
  \]
  This seed-42 interval covers the hidden $\mu=502$. Coverage in one sample
  is not a probability statement about this realized interval.
\end{frame}

\begin{frame}[fragile]{Step 2: The 95\% Z Interval}
  \begin{columns}[T]
    \begin{column}{0.40\textwidth}
      \small
      \begin{lstlisting}
margin = 1.96 * se
lower = xbar - margin
upper = xbar + margin
      \end{lstlisting}

      \begin{block}{Verified output}
        $z^{*}=1.96$\\
        Margin $=3.266667$\\
        CI $[499.471831, 506.005165]$
      \end{block}
    \end{column}
    \begin{column}{0.60\textwidth}
      \centering
      \includegraphics[width=\textwidth]{../figures/z_ci_mean_02_z_interval.png}
      \vspace{0.2em}
      \scriptsize Known-sigma z interval from Step 2
    \end{column}
  \end{columns}
\end{frame}

\begin{frame}{Step 3: Unknown $\sigma$ Makes Z the Wrong Tool}
  \begin{columns}[T]
    \begin{column}{0.42\textwidth}
      \small
      Sample $s=8.390375\neq 10$.

      \vspace{0.5em}
      Invalid z-with-$s$ interval:
      \[
        [499.997642, 505.479354].
      \]

      \vspace{0.4em}
      Width drops from $6.533333$ to $5.481711$.

      \vspace{0.5em}
      \textbf{Limit:} keeping $z^{*}=1.96$ after replacing $\sigma$ by $s$ is
      not a z interval. Lesson 28 uses $t$ with $\mathrm{df}=n-1$.
    \end{column}
    \begin{column}{0.58\textwidth}
      \centering
      \includegraphics[width=\textwidth]{../figures/z_ci_mean_03_unknown_sigma_limit.png}
      \vspace{0.2em}
      \scriptsize Valid z interval versus invalid z-with-$s$ from Step 3
    \end{column}
  \end{columns}
\end{frame}

\begin{frame}{Concept Check}
  \begin{enumerate}
    \item Why is $\mathrm{SE}=10/\sqrt{36}$ and not $s/\sqrt{n}$ in this lesson?
    \item Compute the 95\% margin $1.96\times 1.666667$.
    \item Does covering $\mu=502$ in this sample mean the next sample will
      cover $\mu$?
    \item Why is $s=8.390375$ not a license to keep $z^{*}=1.96$?
  \end{enumerate}

  \vspace{0.6em}
  \begin{block}{Connection to Lesson 28}
    The session can continue directly into the Student $t$ interval, degrees
    of freedom, and the effect of a single outlier on $s$.
  \end{block}

  \vfill
  \scriptsize Source: OpenStax, \textit{Introductory Business Statistics 2e},
  Chapter 8, Section 8.1. CC BY-NC-SA.
\end{frame}

\appendix



\end{document}
